
\documentclass[10pt]{article}
\usepackage[utf8]{inputenc}
\usepackage[T1]{fontenc}
\usepackage{amsmath, amssymb, xcolor, geometry, enumitem, multicol, tcolorbox}

\definecolor{mainblue}{RGB}{0, 102, 204}
\geometry{a4paper, margin=1.2cm, top=1cm, bottom=3cm, footskip=1.2cm}

\begin{document}
\enlargethispage{2.5cm}

% Modern Header
\begin{tcolorbox}[colback=mainblue!5, colframe=mainblue, sharp corners, boxrule=0.5pt]
    \small \textbf{Unit:} undefined \hfill \textbf{Name:} \rule{3cm}{0.4pt} \\
    \small \textbf{Topic:} Variables and Expressions / Order of Operations \hfill \textbf{Date:} \rule{2cm}{0.4pt} \\
    \centering \textbf{\large Advanced (Level 1)}
\end{tcolorbox}

\vspace{0.2cm}

% MCQ Section
\begin{multicols}{2}
\begin{enumerate}[label=\textbf{Q\arabic*.}, leftmargin=*, itemsep=6pt, parsep=0pt]

  \item \small What is the variable in the expression $2x + 3$?
  \begin{enumerate}[label=(\Alph*), leftmargin=0.6cm, nosep, itemsep=2pt]
    \item $2$
    \item $3$
    \item $x$
    \item $2x$
    \item $2 + 3$
  \end{enumerate}

  \item \small Identify the coefficient in the term $4y$
  \begin{enumerate}[label=(\Alph*), leftmargin=0.6cm, nosep, itemsep=2pt]
    \item $y$
    \item $4$
    \item $4y$
    \item $0$
    \item $1$
  \end{enumerate}

  \item \small What is the constant term in the expression $2x - 5$?
  \begin{enumerate}[label=(\Alph*), leftmargin=0.6cm, nosep, itemsep=2pt]
    \item $2x$
    \item $-5$
    \item $2$
    \item $5$
    \item $2 - 5$
  \end{enumerate}

  \item \small Identify the operator in $7 - 2$
  \begin{enumerate}[label=(\Alph*), leftmargin=0.6cm, nosep, itemsep=2pt]
    \item $+$
    \item $-$
    \item $	imes$
    \item $div$
    \item $=$
  \end{enumerate}

  \item \small What is the term $3z$ an example of?
  \begin{enumerate}[label=(\Alph*), leftmargin=0.6cm, nosep, itemsep=2pt]
    \item Constant
    \item Variable
    \item Coefficient
    \item Monomial
    \item Binomial
  \end{enumerate}

  \item \small In the expression $9x^2$, what is $x^2$?
  \begin{enumerate}[label=(\Alph*), leftmargin=0.6cm, nosep, itemsep=2pt]
    \item Variable
    \item Coefficient
    \item Constant
    \item Exponent
    \item Monomial
  \end{enumerate}

  \item \small Identify the exponent in $2x^3$
  \begin{enumerate}[label=(\Alph*), leftmargin=0.6cm, nosep, itemsep=2pt]
    \item $2$
    \item $x$
    \item $3$
    \item $2x$
    \item $x^3$
  \end{enumerate}

  \item \small In $4(2x + 1)$, what operation should be performed first?
  \begin{enumerate}[label=(\Alph*), leftmargin=0.6cm, nosep, itemsep=2pt]
    \item Multiplication
    \item Addition
    \item Subtraction
    \item Division
    \item Exponentiation
  \end{enumerate}

\end{enumerate}
\end{multicols}

\vspace{-0.2cm}
\hrule
\vspace{0.2cm}
\textbf{\small Open Ended Questions (FRQ)}
\begin{enumerate}[label=\textbf{Q\arabic*.}, start=9, itemsep=5pt, topsep=0pt]

  \item \small Draw a diagram to illustrate the concept of variables and constants in the expression $3x + 2$. Label each part. \vspace{2cm}

  \item \small Explain in your own words what the order of operations is and provide an example. Use $2 + 3 	imes 4$ as your example. \vspace{2cm}

\end{enumerate}

\vfill

% Chef's Tips Footer
\begin{tcolorbox}[colback=blue!5, colframe=mainblue!50, title=\small \textbf{Chef's Tips}, fonttitle=\bfseries, bottom=1mm]
\begin{itemize}[noitemsep, topsep=2pt, leftmargin=0.5cm]
  \item \scriptsize Keep the focus on basic definitions and single-step identification for Level 1.\item \scriptsize Ensure students understand that variables represent unknown or changing values, while constants are fixed values.\item \scriptsize Emphasize the importance of following the order of operations (PEMDAS/BODMAS) to simplify expressions correctly.
\end{itemize}
\end{tcolorbox}

\end{document}
  